% Unit 129 English translation of chapter9.tex, lines 1369--1456.
% Source: Methods of Algebra, Volume 2, commit
% 9a5803ff77dd3257484cb177f851a73770a59dd3 (CC BY 4.0).
% stable-unit-id: o014.aljabr2.chapter9.reconstruction-b
% Coverage: the remainder of Section 9.5, abelian subcategories, and the
% reconstruction theorem.

% segment-id: o014.li2.chapter9.reconstruction-b.p001
The relation between the preceding result and abelian subcategories is also
easy to describe.

% segment-id: o014.li2.chapter9.reconstruction-b.l001
\begin{lemma}\label{prop:subcat-gen-L}
	Suppose that $\mathcal{A}$ and $\omega$ satisfy the hypotheses of
	Lemma~\ref{prop:Tannaka-projective-generator}. Every abelian subcategory
	$\mathcal{A}'$, together with $\omega|_{\mathcal{A}'}$, also satisfies
	these hypotheses. Moreover, the canonical homomorphism
	$\coEnd(\omega|_{\mathcal{A}'})\to\coEnd(\omega)$ from
	Remark~\ref{rem:L-cogebra-functoriality} is monic.
\end{lemma}
\begin{proof}
	It is clear that $\mathcal{A}'$ is locally finite. Without loss of
	generality, we may suppose that
	$\Obj(\mathcal{A}')\subset\Obj(\mathcal{A})$ is closed under
	isomorphisms. We use the following general properties; their proofs are
	not difficult and can be found in the exercises of \CHref{sec:app-Abel}.
	\begin{itemize}
		\item Choose a projective generator $s$ of $\mathcal{A}$. It has a
		largest quotient object lying in $\mathcal{A}'$, denoted by $s'$, and
		$s'$ is a projective generator of $\mathcal{A}'$.
		\item Write $A:=\End(s)$ and
		$A':=\End(s')\simeq\Hom(s,s')$. There is a two-sided ideal
		$\mathfrak a$ such that $A\to A'$ induces
		$A/\mathfrak a\rightiso A'$. More explicitly, if
		$t:=\Ker[s\twoheadrightarrow s']\subset s$, then
		$\mathfrak a=\Hom(s,t)\subset A$.
		\item After identifying $\mathcal{A}$ with
		$\catesubd{Mod}{\mathrm{fg}}A$, the category $\mathcal{A}'$ is
		identified with the full subcategory consisting of modules annihilated
		by $\mathfrak a$.
	\end{itemize}

	Choose a basis $\phi_1,\ldots,\phi_m$ of the $\Bbbk$-vector space
	$\Hom(s,t)$. The generator property gives an exact sequence
	% segment-id: o014.li2.chapter9.reconstruction-b.q001
	\[ s^{\oplus m} \xrightarrow{(\phi_i)_i} s \to s' \to 0. \]
	Write $P:=\omega(s)$ and $P':=\omega(s')$. The algebra $A$ acts on $P$
	from the left through $\omega$, and the corresponding exact sequence
	% segment-id: o014.li2.chapter9.reconstruction-b.q002
	\[ P^{\oplus m} \xrightarrow{(\omega(\phi_i))_i} P \to P' \to 0 \]
	shows that $P'\simeq P/\mathfrak aP$.

	The proof of Lemma~\ref{prop:Tannaka-projective-generator} showed that
	$(\mathord\cdot)\dotimes{A}P:\cated{Mod}A\to\cated{Mod}B$ is exact.
	Together with the description above, this gives
	% segment-id: o014.li2.chapter9.reconstruction-b.q003
	\begin{align*}
		(P')^\vee \dotimes{A/\mathfrak{a}} P' & \simeq (A/\mathfrak{a} \dotimes{A} P)^\vee \dotimes{A/\mathfrak{a}} ( A/\mathfrak{a} \dotimes{A} P) \\
		& \simeq (A/\mathfrak{a} \dotimes{A} P)^\vee \dotimes{A} P \\
		& \hookrightarrow P^\vee \dotimes{A} P.
	\end{align*}
	It is not difficult to check that this is precisely the canonical
	homomorphism $\coEnd(\omega|_{\mathcal{A}'})\to\coEnd(\omega)$.
\end{proof}

% segment-id: o014.li2.chapter9.reconstruction-b.eg001
\begin{example}\label{eg:Tannakian-I-graded-0}
	Here are some initial applications of
	Lemma~\ref{prop:Tannaka-projective-generator} for $B=\Bbbk$; recall
	that $\Bbbk$ is assumed to be a field. For a nonempty set $I$, define
	$\cate{Vect}_I(\Bbbk)$ to be the category of $I$-graded
	$\Bbbk$-vector spaces (Example~\ref{eg:graded-module}, with trivial
	$\epsilon$). The objects $(M^i)_{i\in I}$ satisfying
	$\sum_i\dim_{\Bbbk}M^i<\infty$ form the full subcategory
	$\mathcal{A}:=\cate{Vect}_{I,\mathrm f}(\Bbbk)$. Write $\Bbbk[I]$ for
	the $\Bbbk$-vector space with basis $I$; the following structure makes it
	a cocommutative coalgebra:
	% segment-id: o014.li2.chapter9.reconstruction-b.q004
	\[ \epsilon(i) = 1, \quad \Delta(i) = i \otimes i, \quad i \in I. \]
	Define a functor
	$\omega_I:\mathcal{A}\to\cate{Vect}_{\mathrm f}(\Bbbk)$ by sending
	$(M^i)_{i\in I}$ to $\bigoplus_iM^i$.
	\begin{enumerate}
		\item The simplest case $|I|=1$ amounts to taking
		$\mathcal{A}=\cate{Vect}_{\mathrm f}(\Bbbk)$ and
		$\omega_I=\identity$; the projective generator $s$ may be taken to be
		$\Bbbk$. A direct calculation from the definition, or
		Lemma~\ref{prop:Tannaka-projective-generator}(i), gives
		% segment-id: o014.li2.chapter9.reconstruction-b.q005
		\[ \coEnd(\omega_I) \simeq \Bbbk, \quad \epsilon(1) = 1, \quad \Delta(1) = 1 \otimes 1. \]
		Thus, when $|I|=1$, we have
		$\coEnd(\omega_I)\simeq\Bbbk\simeq\Bbbk[I]$.

		\item Next, if $I$ is finite, then
		$\coEnd(\omega_I)\simeq\Bbbk[I]$. This can be deduced from the
		preceding case or understood by applying
		Lemma~\ref{prop:Tannaka-projective-generator}(i), taking as projective
		generator $s$ the object defined by $M^i=\Bbbk$ for every $i$.

		Observe that if $I\subset J$, the embedding
		$\cate{Vect}_{I,\mathrm f}(\Bbbk)\subset
		\cate{Vect}_{J,\mathrm f}(\Bbbk)$ induces
		$\Bbbk[I]\to\Bbbk[J]$, namely the coalgebra embedding induced by
		$I\hookrightarrow J$.

		\item For every nonempty set $I$, we likewise have
		$\coEnd(\omega_I)\simeq\Bbbk[I]$. More precisely, if
		$x\in M^i$ and $\phi\in(M^i)^\vee$ satisfy $\phi(x)=1$, then
		$[\phi\otimes x]\in\coEnd(\omega_I)$ corresponds to
		$i\in\Bbbk[I]$.

		The set $I$ is the filtered union of its finite subsets, and both the
		category $\cate{Vect}_{I,\mathrm f}(\Bbbk)$ and the coalgebra
		$\Bbbk[I]$ reduce to the finite case in the same way. Since the case
		of finite $I$ is known, Proposition~\ref{prop:coHom-colim} immediately
		gives
		% segment-id: o014.li2.chapter9.reconstruction-b.q006
		\[ \coEnd(\omega_I) \simeq \varinjlim_{J \subset I, |J| < \infty} \coEnd(\omega_J). \]
		Taking colimits is a key technique that will reappear in the next
		proof.
	\end{enumerate}
\end{example}

% segment-id: o014.li2.chapter9.reconstruction-b.p002
The following theorem is the main intermediate result so far, and its proof
depends on the technical results of \S\ref{sec:locally-finite-Abel-cat}.
Recall that if $(L,\mu,\eta,\Delta,\epsilon)$ is a $B$-bialgebra,
Proposition~\ref{prop:bialgebra-Mod-monoidal} equips
$\cated{Comod}L$ with a natural monoidal structure: the operation $\otimes$
comes from $\mu$, and the unit is $B$ with the coaction given by
$\eta:B\to B\dotimes{B}L\simeq L$. Observe that
$\catesubd{Comod}{\mathrm{pf}}L$ is a monoidal subcategory.

% segment-id: o014.li2.chapter9.reconstruction-b.th001
\begin{theorem}\label{prop:Tannaka-aux1}
	\index{reconstruction theorem}
	Let $\Bbbk$ be a field, let $B$ be a $\Bbbk$-algebra, let $\mathcal{A}$
	be a locally finite abelian category, and let
	$\omega:\mathcal{A}\to\cated{Mod}B$ be a faithful exact functor taking
	values in $\catesubd{Mod}{\mathrm{pfg}}B$.
	\begin{enumerate}[(i)]
		\item The functor $\overline\omega$ in the canonical factorization
		\eqref{eqn:L-canonical-decomp} is an equivalence of categories.
		Moreover, there is a family of abelian subcategories $\mathcal{A}'$ of
		$\mathcal{A}$ such that
		\begin{compactitem}
			\item $\bigcup_{\mathcal{A}'}\mathcal{A}'=\mathcal{A}$ is a
			filtered union of subcategories;
			\item every $\mathcal{A}'$ has a projective generator,
		\end{compactitem}
		and for this family there is an isomorphism of coalgebras
		% segment-id: o014.li2.chapter9.reconstruction-b.q007
		\[ \coEnd(\omega) \simeq \varinjlim_{\mathcal{A}'} \coEnd(\omega'), \quad \omega' := \omega|_{\mathcal{A}'}. \]
		The transition homomorphisms on the right come from
		Remark~\ref{rem:L-cogebra-functoriality}, and all of them are coalgebra
		embeddings. In this situation,
		% segment-id: o014.li2.chapter9.reconstruction-b.q008
		\[ \catesubd{Comod}{\mathrm f}\coEnd(\omega)=
		\catesubd{Comod}{\mathrm{pf}}\coEnd(\omega). \]

		\item If $\mathcal{A}$ has a monoidal-category structure,
		$B=\Bbbk$, and $\omega$ is a monoidal functor, then with respect to the
		$\Bbbk$-bialgebra structure on $\coEnd(\omega)$ from
		Proposition~\ref{prop:L-bialgebra}, both $\overline\omega$ and $U$ are
		monoidal functors.
		\item Under the same assumptions, if in addition $\mathcal{A}$ is a
		symmetric monoidal category and $\omega$ is compatible with the
		braidings, then $\coEnd(\omega)$ is commutative as an algebra, while
		$\overline\omega$ and $U$ are also compatible with the braidings.
	\end{enumerate}
\end{theorem}
\begin{proof}
	Proposition~\ref{prop:subcat-union} writes $\mathcal{A}$ as the filtered
	union of abelian subcategories $\mathcal{A}'$ having projective
	generators; for example, we may take $\mathcal{A}'=\lrangle{X}$ for
	$X\in\Obj(\mathcal{A})$. Proposition~\ref{prop:coHom-colim} shows that
	$\coEnd(\omega)\simeq\varinjlim_{\mathcal{A}'}\coEnd(\omega')$, and
	Lemma~\ref{prop:subcat-gen-L} shows that all the transition homomorphisms
	are coalgebra embeddings. The canonical factorization
	\eqref{eqn:L-canonical-decomp} is obtained by taking the filtered union
	of
	% segment-id: o014.li2.chapter9.reconstruction-b.q009
	\[\begin{tikzcd}
		\mathcal{A}' \arrow[r, "{\overline{\omega'}}"] & \catesubd{Comod}{\mathrm{pf}}\coEnd(\omega') \arrow[r, "U"] & \catesubd{Mod}{\mathrm{pfg}}B.
	\end{tikzcd}\]
	Lemma~\ref{prop:Tannaka-projective-generator} states that every
	$\overline{\omega'}$ is an equivalence, so $\overline\omega$ is also an
	equivalence. Moreover,
	$\catesubd{Comod}{\mathrm{pf}}\coEnd(\omega')=
	\catesubd{Comod}{\mathrm f}\coEnd(\omega')$; the same process
	simultaneously gives this equality for $\coEnd(\omega)$. This proves (i).

	If $\mathcal{A}$ is monoidal and $B=\Bbbk$, the forgetful functor $U$ is
	always monoidal for the structure on $\cated{Comod}\coEnd(\omega)$ coming
	from the bialgebra $\coEnd(\omega)$. The commutative diagram
	\eqref{eqn:L-bialgebra-prod}, with all of
	$\omega_i,\omega_i',\omega_i''$ replaced by $\omega$, shows that
	$\overline\omega$ is also monoidal. This proves (ii), and a similar
	argument proves (iii).
\end{proof}

% segment-id: o014.li2.chapter9.reconstruction-b.eg002
\begin{example}\label{eg:Tannakian-I-graded}
	Let $\Gamma$ be a monoid and let $B=\Bbbk$ be a field, so that
	$\catesubd{Mod}{\mathrm{pfg}}B=\cate{Vect}_{\mathrm f}(\Bbbk)$.
	Example~\ref{eg:Tannakian-I-graded-0} shows that the coalgebra
	corresponding to
	% segment-id: o014.li2.chapter9.reconstruction-b.q010
	\[ \omega_\Gamma: \cate{Vect}_{\Gamma, \mathrm{f}}(\Bbbk) \to \cate{Vect}_{\mathrm{f}}(\Bbbk), \quad (M^\gamma)_{\gamma \in \Gamma} \to \bigoplus_\gamma M^\gamma \]
	is $\Bbbk[\Gamma]$. Moreover, $\Bbbk[\Gamma]$ carries a bialgebra
	structure (Example~\ref{eg:monoid-bialgebra}). On the other hand,
	$\cate{Vect}_{\Gamma,\mathrm f}(\Bbbk)$ is a monoidal category and
	$\omega_\Gamma$ is a monoidal functor (Example~\ref{eg:graded-module});
	this also equips $\Bbbk[\Gamma]$ with a bialgebra structure. The two
	structures coincide, as one checks by comparing the description of
	multiplication in \eqref{eqn:L-tensor} with
	Example~\ref{eg:Tannakian-I-graded-0}; the details are left to the reader.
\end{example}
